\chapter{Verketten und Koppeln}
\label{b04guide:ch:combine}

\section{Hintereinanderausführen totaler Relationen}

\begin{definition}[Verkettung]
Für $R\subseteq A\times B$ und $S\subseteq B\times C$ sei
\[
 S\circ R
 =\{(a,c)\in A\times C\mid
       \exists b\in B\,((a,b)\in R\land(b,c)\in S)\}.
\]
Die Reihenfolge lautet: zuerst $R$, dann $S$.
Diese Mengenoperation wird hier nur im Hinblick auf Totalität
verwendet. Allgemeine Rechengesetze der Relationsalgebra
sind kein eigener Gegenstand des Leitfadens.
\end{definition}

\begin{theorem}[Abgeschlossenheit unter Verkettung]
\label{b04guide:thm:composition}
Für $\TR(R;A,B)$ und $\TR(S;B,C)$ gilt
\[
 \TR(S\circ R;A,C).
\]
Die Aussage benötigt keine Auswahl von Zwischenelementen
für alle $a\in A$ zugleich.
\end{theorem}

\begin{theorem}[Totalität vor einem totalen zweiten Schritt]
\label{b04guide:thm:compositiontest}
Seien $R\subseteq A\times B$ und $\TR(S;B,C)$. Dann gilt
\[
 \TR(S\circ R;A,C)\quad\Longleftrightarrow\quad\TR(R;A,B).
\]
Ein totaler zweiter Schritt kann fehlende Paare über
Elementen der Ausgangsmenge des ersten Schritts nicht ersetzen.
\end{theorem}

\begin{theorem}[Totalität auf der erreichten Zwischenmenge]
\label{b04guide:thm:reachable}
Seien $\TR(R;A,B)$ und $S\subseteq B\times C$.
Falls
\[
 \forall b\in R[A]\,\exists c\in C\,((b,c)\in S),
\]
so gilt $\TR(S\circ R;A,C)$.
Es genügt also, daß der zweite Schritt auf der von $R$
tatsächlich erreichten Zwischenmenge total ist.
\end{theorem}

\begin{theorem}[Genaues Verkettungskriterium]
\label{b04guide:thm:exactcomposition}
Seien $\TR(R;A,B)$ und $S\subseteq B\times C$. Dann gilt
\[
 \begin{gathered}
 \TR(S\circ R;A,C)\\
 \Longleftrightarrow\
 \forall a\in A\,\exists b\in R_a\,
                  \exists c\in C\,((b,c)\in S).
 \end{gathered}
\]
Jede Faser des ersten Schritts muß mindestens ein
Zwischenelement enthalten, von dem aus der zweite Schritt
fortgesetzt werden kann.
\end{theorem}

\section{Zwei Zuordnungen gleichzeitig}

\begin{definition}[Kopplung bei gemeinsamer Ausgangsmenge]
Für $R\subseteq A\times B$ und $S\subseteq A\times C$ sei
\[
 K(R,S)=
 \{(a,(b,c))\in A\times(B\times C)\mid
                    (a,b)\in R\land(a,c)\in S\}.
\]
Die beiden Ziele werden als ein geordnetes Paar zusammengefaßt.
\end{definition}

\begin{theorem}[Kopplung und Totalität]
\label{b04guide:thm:coupling}
Für $R\subseteq A\times B$ und $S\subseteq A\times C$ gilt
\[
 \begin{gathered}
 \TR\bigl(K(R,S);A,B\times C\bigr)\\
 \Longleftrightarrow\
             \TR(R;A,B)\land\TR(S;A,C).
 \end{gathered}
\]
Insbesondere liefern zwei totale Relationen mit derselben
Ausgangsmenge wieder eine totale Relation.
\end{theorem}

\begin{definition}[Paarweise Produktrelation]
Für $R\subseteq A\times B$ und $S\subseteq C\times D$ sei
\[
 R\boxtimes S=
 \{((a,c),(b,d))\in(A\times C)\times(B\times D)
             \mid(a,b)\in R\land(c,d)\in S\}.
\]
$\boxtimes$ bezeichnet hier nur diese ausdrücklich definierte
Menge von Paaren.
\end{definition}

\begin{theorem}[Produkt totaler Relationen]
\label{b04guide:thm:product}
Für $\TR(R;A,B)$ und $\TR(S;C,D)$ gilt
\[
 \TR(R\boxtimes S;A\times C,B\times D).
\]
Sind umgekehrt $A$ und $C$ nichtleer, so gilt für beliebige
$R\subseteq A\times B$ und $S\subseteq C\times D$ auch
\[
 \begin{gathered}
 \TR(R\boxtimes S;A\times C,B\times D)\\
 \Longrightarrow\quad\TR(R;A,B)\land\TR(S;C,D).
 \end{gathered}
\]
Bei leerem $A$ oder leerem $C$ ist die Produktrelation stets
total auf ihrer leeren Ausgangsmenge, unabhängig von der
Totalität der einzelnen Relationen.
\end{theorem}

Die beiden Produktkonstruktionen verbinden jeweils zwei
Zuordnungen. Beliebige gleichzeitige Auswahlvorgänge und
unendliche Fortsetzungen sind hier nicht Gegenstand.
